\section{Introduction}
\label{sec:intro}

Retrieval-augmented generation (RAG)~\cite{lewis2020rag} grounds a language model in retrieved passages, but many questions cannot be answered from any single passage. In \emph{bridge} questions~\cite{yang2018hotpotqa} the second passage is reachable only through an entity named in the first, and the question itself never mentions that entity: ``Which university did the director of \emph{Film X} attend?'' needs the passage about the film, which names the director, \emph{and} the passage about the director. Similarity-based retrieval ranks the second passage low because it shares little vocabulary with the question. The generator then sees the endpoints of the reasoning chain without the link between them and either abstains or guesses. We call this \emph{bridge failure}.

Remedies fall into two families. \emph{Iterative retrieval} lets an LLM reformulate the query after each round~\cite{trivedi2023ircot,asai2024selfrag}; it works but costs several LLM calls per question. \emph{Graph-based RAG} builds entity or passage graphs and uses them to find or re-rank candidates~\cite{edge2024graphrag,guo2024lightrag,gutierrez2025hipporag2,he2024gretriever,rappazzo2024gem}; the graph typically scores passages \emph{one at a time}. Neither family reasons directly about the property that matters for a bridge question: whether the \emph{set} of passages that will be shown to the generator is connected through the bridge.

We ask four questions. \textbf{RQ1:} how much of the answer-quality loss on bridge questions is attributable to a bridge-failed context, as opposed to generator behaviour? \textbf{RQ2:} can a cheap selector that uses no LLM call, and rewards a \emph{connected} context under a token budget, repair bridge failure better than flat retrieval and graph-based baselines? \textbf{RQ3:} is spectral connectivity (the Fiedler value $\lambda_2$ of the selected subgraph) the active ingredient, or would any connectivity reward do? \textbf{RQ4:} when does the effect transfer across graphs, generators and datasets?

We propose SES-RAG (Spectral Evidence Selection RAG), a greedy context selector that adds to the relevance of each candidate the change in $\lambda_2$ of the induced evidence subgraph, and evaluate it against seven baselines, a gold-first oracle and four generators from two vendors, under a protocol with a held-out primary set, equal tuning effort, Holm-corrected paired tests and a pre-declared pass criterion. Our contributions are:
\begin{itemize}
\item \textbf{A measurement of bridge failure} (RQ1). Conditioning on whether the context contains all gold passages separates retrieval from generation: with a complete context the four generators reach 72.6--79.9\% F1, with a bridge-failed one 19.2--24.1\%, and they abstain on 53--63\% of failed contexts.
\item \textbf{SES-RAG}, a set-level, LLM-free selector (RQ2). It is the best method that makes no LLM call on HotpotQA bridge questions, improving gold-complete retrieval by about 10 points over Hybrid RAG and F1 by 4.0--6.5 points, with confidence intervals excluding zero for all four generators. It does not significantly beat the strongest graph baselines, and an LLM-driven iterative baseline is ahead of it.
\item \textbf{Negative and boundary results} (RQ3, RQ4). Component count and raw-$\lambda_2$ connectivity terms work as well as or better than normalised $\lambda_2$, so spectral connectivity is not special; the advantage appears on 2WikiMultiHopQA but not on MuSiQue, disappears when the graph lacks title-mention edges, and an extension that packs iterative-retrieval candidates with the same rule gives no significant gain on the primary set.
\item \textbf{A pre-specified, reproducible evaluation protocol} with a held-out primary set, equal tuning effort per system, frozen parameters, and code and raw outputs released.
\end{itemize}
